\chapter{Programs — Expressions}
%%% generated by the blueprint pipeline from TCSlib.Programs.Light.Compiler.Expressions
\section{Overview}

Declarations of the Lean module \texttt{TCSlib.Programs.Light.Compiler.Expressions}.

\section{Declarations}

\begin{theorem}[Positivity of expression height]
\lean{Light.Expr.height_pos}
\label{Light.Expr.height_pos}
\leanok
\uses{Light.Expr, Light.Expr.height}
\difficulty{4}
Every expression $e$ has strictly positive height: $0 < \mathrm{height}(e)$.
\end{theorem}

\begin{theorem}[Expression temporary belongs to the temporary set]
\lean{Light.Compiler.cT_mem_temps_of_height}
\label{Light.Compiler.cT_mem_temps_of_height}
\leanok
\uses{Light.Compiler.Temps, Light.Compiler.cT, Light.Expr, Light.Expr.height, Light.Expr.height_pos}
\difficulty{1}
Let $F,k\in\bbn$ and let $e$ be an expression built from natural-number constants and
variables, binary addition, subtraction, or multiplication, and loads. If
$k+\mathrm{height}(e)\le 2F+1$, then the temporary address $-(23+4k)$ belongs to the set
consisting of the address $-13$ and all addresses $-(23+4j)$ with $k\le j\le 2F$.
\end{theorem}

\begin{theorem}[Straight instructions in digit code]
\lean{Light.Compiler.straight_digitsCode}
\label{Light.Compiler.straight_digitsCode}
\leanok
\uses{Light.Compiler.digitsCode, WordRAM.Straight}
\difficulty{4}
For every integer location $c$ and every finite list of natural-number digits, each
instruction in the code generated from that digit list at location $c$ is straight.
\end{theorem}

\begin{theorem}[Effects of digit code]
\lean{Light.Compiler.effects_digitsCode}
\label{Light.Compiler.effects_digitsCode}
\leanok
\uses{Light.Compiler.cONE, Light.Compiler.digitsCode, WordRAM.AgreeOutside, WordRAM.AgreeOutside.read, WordRAM.AgreeOutside.refl, WordRAM.AgreeOutside.update, WordRAM.effect_add, WordRAM.effect_sub, WordRAM.effects, WordRAM.effects_append, WordRAM.effects_cons, WordRAM.effects_nil, WordRAM.wd, WordRAM.wd_toInt}
\difficulty{6}
Let $c$ be a memory address distinct from the address holding the constant $1$, and let
$L$ be a list of binary digits, each less than $2$. If the initial memory holds the word
value $1$ at the constant-one address, executing the digit code for $L$ at address $c$
leaves the natural number represented by those digits at $c$. It leaves memory unchanged
outside $\{c\}$.
\end{theorem}

\begin{theorem}[Straightness of arithmetic instructions]
\lean{Light.Compiler.straight_instr}
\label{Light.Compiler.straight_instr}
\leanok
\uses{Light.Op, Light.Op.instr, WordRAM.Straight}
\difficulty{4}
For every operation chosen from addition, subtraction, and multiplication, and every
three integer locations $i,j,k$, the corresponding instruction with destination $i$ and
operands $j,k$ is straight.
\end{theorem}

\begin{theorem}[Straight instructions in expression code]
\lean{Light.Compiler.straight_compileExpr}
\label{Light.Compiler.straight_compileExpr}
\leanok
\uses{Light.Compiler.compileExpr, Light.Compiler.straight_digitsCode, Light.Compiler.straight_instr, Light.Expr, WordRAM.Straight}
\difficulty{4}
For every expression and every natural number $k$, each instruction in the code that
compiles the expression starting at $k$ is straight.
\end{theorem}

\begin{theorem}[Effect of a compiled arithmetic instruction]
\lean{Light.Compiler.effect_instr}
\label{Light.Compiler.effect_instr}
\leanok
\uses{Light.Op, Light.Op.eval, Light.Op.instr, WordRAM.effect, WordRAM.effect_add, WordRAM.effect_mul, WordRAM.effect_sub, WordRAM.wd}
\difficulty{4}
For an arithmetic operation $o$ and word-RAM memory $m$, if the values in cells $j$ and
$k$ are the $W$-bit encodings of integers $a$ and $b$, respectively, then executing the
instruction representing $o$ with destination cell $i$ changes the memory exactly by
writing the $W$-bit encoding of $o(a,b)$ to cell $i$.
\end{theorem}

\begin{theorem}[Compiled expression effects]
\lean{Light.Compiler.effects_compileExpr}
\label{Light.Compiler.effects_compileExpr}
\leanok
\uses{Light.Compiler.Framed, Light.Compiler.Framed.same, Light.Compiler.Rel.read, Light.Compiler.Sizes.inRange_addr, Light.Compiler.Sizes.inRange_stack, Light.Compiler.Temps, Light.Compiler.cADDR, Light.Compiler.cStack_sub, Light.Compiler.cT, Light.Compiler.cT_mem_temps_of_height, Light.Compiler.compileExpr, Light.Compiler.effect_instr, Light.Compiler.effects_digitsCode, Light.Compiler.temps_subset_scratch, Light.Expr, Light.Expr.Safe, Light.Expr.height, Light.Expr.val, Light.Expr.vars, Light.Expr.vars_op_le_iff, WordRAM.AgreeOutside, WordRAM.AgreeOutside.cell, WordRAM.AgreeOutside.mono, WordRAM.AgreeOutside.read, WordRAM.AgreeOutside.refl, WordRAM.AgreeOutside.trans, WordRAM.AgreeOutside.update, WordRAM.effect_load, WordRAM.effect_sub, WordRAM.effects, WordRAM.effects_append, WordRAM.effects_cons, WordRAM.effects_nil, WordRAM.wd}
\difficulty{6}
Let a word size, compiler size parameters, state, stack position, and memory satisfy the
framed-state conditions: the designated constants, pool, frame pointer, local-variable
cells, and valid memory cells encode the given state, and the stack position is within
bounds. For an expression $e$ and temporary index $k$, suppose $e$ uses at most $F$
variables, is safe under the state and limits, and $k + \mathrm{height}(e) \le 2F+1$.
Executing the compiled expression writes its value into temporary cell $k$ and leaves
memory unchanged outside the set consisting of the address scratch cell and temporary
cells indexed from $k$ through $2F$.
\end{theorem}

\begin{theorem}[Compiled expression effects from width]
\lean{Light.Compiler.effects_compileExpr_of_width}
\label{Light.Compiler.effects_compileExpr_of_width}
\leanok
\uses{Light.Compiler.Framed, Light.Compiler.Temps, Light.Compiler.cT, Light.Compiler.compileExpr, Light.Compiler.effects_compileExpr, Light.Expr, Light.Expr.Safe, Light.Expr.val, Light.Expr.width, Light.Expr.width_le_iff, WordRAM.AgreeOutside, WordRAM.effects, WordRAM.wd}
\difficulty{3}
Let a word size, compiler size parameters, state, stack position, and memory satisfy the
framed-state conditions. For an expression $e$ and temporary index $k$, suppose
$e.\mathrm{width} \le F$, $e$ is safe under the state and limits, and $k \le F+1$.
Executing the compiled expression writes its value into temporary cell $k$ and leaves
memory unchanged outside the set consisting of the address scratch cell and temporary
cells indexed from $k$ through $2F$.
\end{theorem}
